\section{Implementation}
\label{sec:implementation}

\subsection{RTL Organization}

The HyNoC RTL is written in Verilog and organized as follows:

\begin{itemize}
  \item \texttt{hynoc\_ingress}: ingress port (dual-clock FIFO, FSM, routing
        protocol decoder for unicast and multicast);
  \item \texttt{hynoc\_egress}: egress port (data multiplexer, PRRA);
  \item \texttt{hynoc\_router\_base}: parameterized router base (crossbar wiring,
        port instantiation);
  \item \texttt{hynoc\_router\_3p} / \texttt{hynoc\_router\_5p}: 3-port and
        5-port router instances with self-checking testbenches;
  \item \texttt{hynoc\_local\_interface}: local interface for node attachment,
        including the egress output FIFO for clock domain decoupling;
  \item \texttt{hynoc\_stream\_reader} / \texttt{hynoc\_stream\_writer}:
        DMA-style packet injection and reception modules.
\end{itemize}

The design is parameterized by: payload width (\texttt{PAYLOAD\_WIDTH}), FIFO
depth (\texttt{LOG2\_FIFO\_DEPTH}: 2 to 64 entries), number of ports
(\texttt{NB\_PORTS}: 3 to 9), index width (\texttt{INDEX\_WIDTH}), optional
multicast routing (\texttt{ENABLE\_MCAST\_ROUTING}), and single-clock mode
(\texttt{SINGLE\_CLOCK\_ROUTER}).

\subsection{Simulation and Verification}

Each module is delivered with a self-checking testbench. The 3-port router
testbench exercises both unicast and multicast scenarios over a 2-router, 4-node
topology. The 5-port router testbenches cover five distinct traffic scenarios,
including simultaneous multi-path transfers that verify the PRRA's starvation-free
behavior.

Simulation is supported under Icarus Verilog, ModelSim, and Verilator. Waveform
inspection is facilitated by Python-based waveform configuration scripts.

\subsection{FPGA Target}

The design targets Intel (Altera) and AMD (Xilinx) FPGA families. Reference board
support is included for the DE0-Nano (Cyclone~IV) and the ZedBoard (Zynq-7000).
Dual-clock FIFOs are implemented using the \texttt{dclkfifolut} module, which
maps to distributed LUT RAM, avoiding dependency on vendor-specific FIFO IP and
enabling portability across FPGA families.

\subsection{License}

HyNoC is released as open-source hardware under the
\textbf{CERN-OHL-P~v2}~license~\cite{cernohl}, a permissive open hardware license.
The complete source, including RTL, testbenches, documentation, and build
infrastructure, is available at:
\url{https://github.com/cclienti/verilog-ip}.
